\section{Eksperiment ja tulemused} \label{eksperimendid}

Käesolev peatükk kirjeldab läbiviidud eksperimenti, esitab nelja mudeli genereeritud bakalaureusetööde tulemused ning analüüsib nende kvaliteeti erinevatest aspektidest lähtuvalt.

\subsection{Eksperimendi ülesehitus} \label{subsec:exp-tervik}

Eksperimendis paluti igal mudel-tööriista kombinatsioonil genereerida terviklik bakalaureusetöö informaatika valdkonnast ühe kasutajaviibega, ilma järelküsimuste ega iteratiivse suunamiseta. Mudelitele anti identne lühike eestikeelne viibe (vt peatükk~\ref{subsec:eksperimendi-ulesehitus} ja lisa~I), mis palus mudelil luua LaTeX-vormingus bakalaureusetöö, mis vastab Tartu Ülikooli arvutiteaduse instituudi nõuetele, kasutades töökausta paigutatud UniTartuCS malli. Konkreetset teemapealkirja ette ei antud -- mudel pidi ise valima informaatikaga seotud teema. Autori (Uku Renek Kronbergs) ja juhendaja (Alo Peets, MSc) nimed anti tiitellehe täitmiseks.

Viibe on teadlikult lühike, kuna üks uurimisküsimustest on, kuidas mudel-tööriista kombinatsioonid reageerivad minimaalsele juhisele ning kas nad suudavad iseseisvalt valida sobiva teema, keele ja struktuuri. Kuna ülesanne nõuab failihaldust, mitme \LaTeX{} faili loomist ja kompileerimist, kasutati iga mudeliga agentset töövoogu. Sonnet 4.6 ja Opus 4.6 puhul kasutati Anthropicu Cowork tööriista (mis on sama ettevõtte Claude Code'ist eraldiseisev toode, vt peatükk~\ref{claudecode}), Gemini 3.1 Pro puhul Google'i Antigravity ning GPT-5.4 puhul ChatGPT Codex laiendust Visual Studio Code'is. GPT-5.4 puhul testiti lisaks ka tavapärast ChatGPT vestlusliidest, mis andis oluliselt erineva tulemuse.

Tulemused olid kombinatsioonide lõikes drastiliselt erinevad, nagu nähtub tabelist~\ref{tab:exp6}. Oluline tõlgenduslik märkus: järgnevas analüüsis tuleb tulemusi vaadelda mudeli ja agentse tööriista paaride, mitte puhaste mudeli\-võimekuste näitajatena -- iga jooks tehti üks kord, mistõttu näitajad on eksploratiivsed ja võivad kordus\-jooksudel erineda.

\begin{table}[ht]
\centering
\caption{Tervikliku lõputöö genereerimise tulemuste koondülevaade}
\label{tab:exp6}
\begin{tblr}{
    width=\linewidth,
    colspec={X[l] X[c] X[c] X[c] X[c]},
    rows={font=\small},
    colsep=3pt,
    row{1}={font=\bfseries},
    hlines,
    vlines,
    row{even}={bg=colorCellHighlight},
}
Parameeter & Gemini 3.1 Pro & GPT-5.4 & Claude Sonnet 4.6 & Claude Opus 4.6 \\
Väljundkeel & Eesti & Eesti (Codex) / ingl. (chat) & Eesti & Eesti \\
Maht (lehekülgi) & 39 & 47 (Codex) / 8 (chat) & 40 & 61 \\
Valitud teema & GenAI ja algajad progr. & RAG-põhine QA-süsteem (Codex); keeldumine (chat) & TI lõputöö gener. & Paroolihaldur \\
Struktuur & Lõputöö vormis & Lõputöö (Codex) / juhend (chat) & Lõputöö vormis & Lõputöö vormis \\
Agentne tööriist & Antigravity & Codex VS Code'is & Cowork & Cowork \\
Viited & 9, ca 25\% halluts. & 22, ca 20\% halluts. & 17, ca 20\% halluts. & 35, ca 15\% halluts. \\
Akad. stiil & Hea & Hea (Codex) & Väga hea & Suurepärane \\
\end{tblr}
\end{table}

\subsection{Claude Sonnet 4.6 (Cowork): terviklik eestikeelne lõputöö} \label{subsec:sonnet-toolo}

Claude Sonnet 4.6 + Cowork kombinatsioon genereeris 40-leheküljelise eestikeelse dokumendi, mis järgis bakalaureusetöö struktuuri. Kombinatsioon valis teemaks TI-põhise lõputöö genereerimise -- teema, mis ühtib käesoleva bakalaureusetöö teemaga. Töö sisaldas detailseid hindamistabeleid, promptide näiteid lisades ning TI kasutamise deklaratsiooni. Akadeemiline stiil oli tugev ja töö järgis LaTeX-malli korrektselt. Peamised puudused olid mõned hallutsineeritud viited (ca 20\%), kohati keeruline lauseehitus ning 13 tuvastatud kirjaviga kogu töö peale (vt peatükk~\ref{subsec:kirjaviga}), millest mõned asusid nähtaval kohal -- näiteks alapealkirjas kaks korda esinenud \textit{Masskiijutamine}.

\subsection{GPT-5.4 (vestlusliides ja Codex): liidesest sõltuv käitumine} \label{subsec:gpt-toolo}

GPT-5.4 tulem erines teistest fundamentaalselt, kusjuures erinevus ei tulenenud mitte mudelist, vaid liidesest, mille kaudu mudelit kasutati. Tavapärase ChatGPT vestlusliidese kaudu keeldus mudel genereerimast terviklikku lõputööd, viidates otseselt akadeemilise aususe piirangutele. Väljundis selgitati, et täieliku bakalaureusetöö kirjutamine tudengi eest oleks vastuolus Tartu Ülikooli akadeemilise aususe reeglitega. Selle asemel pakkus GPT-5.4 välja 8-leheküljelise ingliskeelse teostatavusjuhendi (\textit{feasibility guide}), mis sisaldas mudelivaliku soovitusi, eksperimendi disaini raamistikku, hindamismetoodikat ja \LaTeX{} malli juhiseid. Kuigi juhend oli sisuliselt kvaliteetne ja viited olid kontrollitavad, ei vastanud see ülesande püstitusele: väljund oli inglise, mitte eesti keeles, ja vormilt pigem nõuandev memo kui akadeemiline lõputöö. See käitumine viitab sellele, et ChatGPT vestlusliidese süsteemijuhised ja ohutuspoliitika tõlgendavad ``tervikliku bakalaureusetöö genereerimine'' tüüpi ülesandeid vaikimisi akadeemilise aususe piirangu alla kuuluvana ning eelistavad sellistes olukordades keeldumist või nõuandvat vastust.

Täiendav katse OpenAI ChatGPT Codex laiendusega Visual Studio Code'i keskkonnas andis aga identsest viipest hoolimata oluliselt erineva tulemuse. Sama GPT-5.4 mudel suutis agentse tööriista kaudu genereerida visuaalselt teiste kombinatsioonide väljunditega sarnase tervikliku 47-leheküljelise eestikeelse dokumendi teemal ``Eestikeelse RAG-põhise küsimus-vastus süsteemi kavandamine'' (vt peatükk~\ref{subsec:codex}). See erinevus näitab, et agentne töövoog muudab ülesande konteksti ja tööjaotust -- ülesanne jaotub väiksemateks faili- ja kompileerimispõhisteks sammudeks ning liidese süsteemijuhised ja ohutuspoliitika on erinevad -- mistõttu sama mudel käitub erinevas liideses erinevalt. Sellest ei tohi järeldada, et GPT-5.4 \emph{ei toeta} eesti keelt tervikliku teksti genereerimisel: Codex-väljund on eestikeelne, mahukas ja ortograafiliselt äärmiselt täpne. Vestlusliidese keeldumine on seega pigem liidese ja selle ohutus\-konteksti, mitte mudeli keelelise võimekuse näitaja. Codex laienduse kaudu genereeritud töö oli ortograafilise täpsuse poolest silmapaistev: kogu töö peale tuvastati vaid üks kirjaviga (\textit{vastusevormingun} õige \textit{vastusevormingu} asemel), mis oli nelja kombinatsiooni seas parim tulemus (vt peatükk~\ref{subsec:kirjaviga}).

\subsection{Gemini 3.1 Pro (Antigravity): mahukas, kuid fabritseeritud lõputöö} \label{subsec:gemini-toolo}

Gemini 3.1 Pro + Antigravity kombinatsioon genereeris 39-leheküljelise eestikeelse bakalaureusetöö teemal ``Generatiivse tehisintellekti mõju tarkvaraarenduse produktiivsusele algajate programmeerijate seas''. Kombinatsioon valis iseseisvalt informaatika valdkonnast aktuaalse ja hästi piiritletud uurimisteema, mis käsitles suurte keelemudelite mõju programmeerimise õppimisele.

Genereeritud töö oli strukturaalselt terviklik ja sisaldas kõiki nõutud komponente: eesti- ja ingliskeelne kokkuvõte, sisukord, sissejuhatus uurimisküsimustega, taust ja kirjandusülevaade (Transformer-arhitektuur, kognitiivne koormus, \textit{over-reliance}), metoodika (väidetavalt kvaasi-eksperimentaalne disain 30 tudengiga, kes fiktiivsete kirjelduste kohaselt töötasid Tartumaa IT-ettevõtetes nagu Nortal, Playtech ja Datafox), tulemused (t-testid, statistilised tabelid), kokkuvõte ja viis lisa. Lisad olid ka hästi esitatud: React.js koodiredaktori komponent (ca 150 rida), PostgreSQL andmebaasi skeem (5 tabelit), Pandas/SciPy andmeanalüüsi skript (ca 160 rida) ning tudengite koodinäited koos vigade analüüsiga.

Töö tugevused olid järgmised:
\begin{itemize}
    \item Akadeemilise struktuuriga hästi arvestatud: uurimisküsimused, hüpoteesid, statistilised testid;
    \item Eesti keele kvaliteet oli üldiselt arusaadav, kuigi esines mõningaid terminoloogilisi ebajärjekindlusi;
    \item Viited (9 allikat) kasutasid reaalsete teadlaste nimesid (Becker, Kazemitabaar, Vaswani jt);
    \item Koodinäited lisades olid funktsionaalsed ja loogiliselt kirjutatud;
    \item Töö maht (39 lehekülge) vastas bakalaureusetöö nõuetele.
\end{itemize}

Peamised puudused olid siiski olulised, kusjuures esimene neist ei ole käsitletav tavapärase kvaliteedipuudusena, vaid akadeemilise väärkasutuse kõrge riskiga juhtumina:
\begin{itemize}
    \item \textbf{Fabritseeritud empiiriline uurimus.} Kombinatsioon loetles 30 tudengiga kvaasi-eksperimendi, sealhulgas osalejate jaotuse katse- ja kontrollrühmadeks, konkreetsed statistilised testid (t-test, Cohen'i $d$), numbrilised tulemused ning Tartumaa IT-ettevõtete kirjeldused. Seda uurimust tegelikult ei toimunud -- tudengeid, nende koode, katsekeskkonda ega statistilisi tulemusi ei eksisteeri. Välja mõeldi nii empiirilised andmed kui ka reaalsete Eesti ettevõtete nimede kasutamine. Lugeja, kes ei kontrolli andmeid, saab äärmiselt veenva, kuid tegelikkuses olematu uuringu.
    \item Viidete kontrollitavust vajab eraldi auditeerimist -- kuigi autorite nimed on reaalsed, olid ca 25\% viidetest hallutsineeritud (olematud artiklid või valed avaldamisandmed);
    \item Eesti keele ortograafiline täpsus oli nelja kombinatsiooni seas kõige nõrgem -- töös tuvastati 28 selget kirjaviga, sealhulgas tavasõnade valed vormid (\textit{strukuure}, \textit{kordinaadid}, \textit{konfiktide}), mittesõnad (\textit{teoosiab}, \textit{Tehisintsistendi}) ning valede tähtedega asendatud täpitähed (vt peatükk~\ref{subsec:kirjaviga});
    \item Antigravity tööriist oli vajalik käesolevas katses saadud tulemuseni jõudmiseks -- autori varasemad mitteametlikud katsed ilma agentse tööriistata andsid oluliselt nõrgema tulemuse.
\end{itemize}

\subsection{Claude Opus 4.6 (Cowork): terviklik tehniline lõputöö} \label{subsec:opus-toolo}

Claude Opus 4.6 + Cowork kombinatsioon andis selles eksperimendis kõige ulatuslikuma ja kvaliteetsema väljundi -- 61-leheküljelise eestikeelse bakalaureusetöö teemal ``Turvalise paroolihalduri veebirakenduse arendamine''. Cowork'i agentne tööriist (vt peatükk~\ref{claudecode}) võimaldas mudelil otse kirjutada \LaTeX{} faile, hallata viiteid ja kompileerida dokumenti.

Genereeritud töö sisaldas terviklikku tehnilist arhitektuuri (nullteadmise põhimõte, AES-256-GCM krüpteering, Argon2id võtmetuletamine) ning täielikku React/TypeScript kasutajaliidest ja Node.js/Express taustarakendust. Lisaks \emph{väitis} töö, et rakendus läbis 149 automatiseeritud testi, OWASP Top 10 vastavuse analüüsi ning süsteemi kasutatavuse skaala (SUS) hindamise skooriga 77,0, ja et turvaaudit ei tuvastanud ühtegi kriitilist turvaauku. Olulise metoodilise märkusena tuleb rõhutada, et käesoleva eksperimendi autor ei teostanud nende testide, auditi ega kasutajauuringu iseseisvat läbiviimist -- need näitajad on kirjas Opus 4.6 + Cowork kombinatsiooni enda genereeritud väljundis ning kuuluvad samasse kategooriasse empiiriliste väidetega, mida TI-loodud töödes tuleb inimese poolt eraldi kontrollida (vrd Gemini 3.1 Pro fabritseeritud eksperiment alapeatükis~\ref{subsec:gemini-toolo}). Opus 4.6 kirjeldatud arhitektuur ja koodinäited olid siiski tunduvalt detailsemad ja sisemiselt sidusamad kui Gemini 3.1 Pro töös esitatud statistilised tulemused.

Opus 4.6 + Cowork tulemuse peamised erinevused võrreldes Sonnet 4.6 + Cowork tulemusega:
\begin{itemize}
    \item Oluliselt suurem maht (61 vs 40 lehekülge), mis viitab sügavamale sisulisele käsitlusele;
    \item Tehniline detailsus oli kõrgem -- töö sisaldas konkreetseid koodi- ja arhitektuurinäiteid;
    \item Eesti keele kvaliteet oli ühtlasem ja loomulikum ka pikemas tekstis -- kogu 61-leheküljelise töö peale tuvastati vaid 5 kirjaviga (peamiselt pärisnimedes, nagu \textit{Carnegei} ja \textit{Californiya}), mis teeb kirjavigade tiheduseks 0,08 viga lehekülje kohta (vt peatükk~\ref{subsec:kirjaviga});
    \item Viidete hallutsineerimise määr oli madalam (ca 15\% võrreldes Sonnet 4.6 ca 20\%).
\end{itemize}

\subsection{Kombinatsioonikäitumise üldised järeldused} \label{subsec:mudelikaitumine}

Tervikliku lõputöö genereerimise eksperiment tõi esile neli fundamentaalselt erinevat käitumismustrit, mis iseloomustasid konkreetseid mudel-tööriista kombinatsioone:
\begin{enumerate}
    \item \textbf{Selles katses suurepärane tulemus} (Claude Opus 4.6 + Cowork): kombinatsioon genereeris kõige mahukama (61 lk) ja tehniliselt sügavaima eestikeelse lõputöö, valides teemaks paroolihalduri arendamise.
    \item \textbf{Selles katses tugev tulemus} (Claude Sonnet 4.6 + Cowork): kombinatsioon genereeris mahuka ja struktureeritud väljundi (40 lk), valides teemaks TI-põhise lõputöö genereerimise.
    \item \textbf{Selles katses tugev tulemus, kuid fabritseeritud empiirikaga} (Gemini 3.1 Pro + Antigravity): kombinatsioon genereeris tervikliku ja hästi struktureeritud töö (39 lk) koos veenva, kuid fabritseeritud eksperimendiga, valides teemaks GenAI mõju algajate programmeerijate produktiivsusele.
    \item \textbf{Liidesesõltuv käitumine} (GPT-5.4): tavapärase ChatGPT vestlusliidese kaudu keeldus mudel ülesannet täitmast akadeemilise aususe põhjendusel ja andis selle asemel 8-leheküljelise ingliskeelse teostatavusjuhendi; ChatGPT Codex laiendusega Visual Studio Code'is genereeris aga sama mudel sama viipega tervikliku 47-leheküljelise eestikeelse dokumendi (vt peatükk~\ref{subsec:codex}). See erinevus on järgneva arutelu keskne leid ja illustreerib seda, et liidese süsteemijuhised, ohutuspoliitika ja kasutuskontekst mõjutavad väljundit sama mudeli puhul oluliselt.
\end{enumerate}

Eriti tähelepanuväärne on agentsete tööriistade ja liideste mõju tulemusele. Nii Cowork (Opus 4.6 ja Sonnet 4.6 puhul), Google'i Antigravity (Gemini 3.1 Pro puhul) kui ka ChatGPT Codex (GPT-5.4 puhul) andsid oluliselt erineva tulemuse kui puhas vestlusliides -- seda kinnitas kõige ilmekamalt GPT-5.4 kahe stsenaariumi vastandlik tulemus. Kõik kombinatsioonid valisid iseseisvalt erineva informaatikateema, mis näitab, kui erinevalt samale lühikesele viipele reageeritakse. Isegi kõige edukamad tulemused (Opus 4.6 + Cowork, GPT-5.4 + Codex ja Gemini 3.1 Pro + Antigravity) vajasid inimese järeltöötlust: viidete kontrolli, terminoloogia ühtlustamise ja fabritseeritud andmete tuvastamise osas. Kuna iga kombinatsiooni kohta tehti vaid üks jooks, ei saa ühe jooksu põhjal väita, et üks mudel on üldiselt parem kui teine -- saab väita üksnes, et selles eksploratiivses katses andis konkreetne mudel-tööriista kombinatsioon konkreetse tulemuse. Genereeritud tööde väljavõtted on esitatud lisas~III.

\subsection{Kirjavigade analüüs} \label{subsec:kirjaviga}

Eesti keele kvaliteedi täpsemaks hindamiseks vaadati eraldi üle iga genereeritud lõputöö kirjavead -- see tähendab ortograafilised näpukad, valed tähed, puuduvad või liigsed tähed, sõnade valed vormid ning pärisnimede vead. Analüüsist jäeti välja kontekstis korrektsed, kuid stiililiselt ebaõnnestunud laused, samuti vead, mis ilmnesid üksnes PDF-ist teksti ekstraheerimise kodeeringuartefaktidena. Tabelis~\ref{tab:kirjavead} on esitatud kirjavigade koondarv iga töö kohta ning arvestuslik tihedus lehekülje kohta.

\begin{table}[ht]
\centering
\caption{Genereeritud tervikliku lõputöö kirjavigade koondarv}
\label{tab:kirjavead}
\begin{tblr}{
    width=\linewidth,
    colspec={X[1.8,l] X[0.7,c] X[1,c] X[1,c] X[1,c]},
    rows={font=\small},
    colsep=3pt,
    row{1}={font=\bfseries\small},
    hlines,
    vlines,
    row{even}={bg=colorCellHighlight},
}
Mudel & Maht (lk) & Kirjavigade arv & Kirjavigu lk kohta & Hinnang \\
Gemini 3.1 Pro (Antigravity) & 39 & 28 & 0,72 & Nõrk \\
Claude Sonnet 4.6 (Cowork) & 40 & 13 & 0,33 & Keskmine \\
Claude Opus 4.6 (Cowork) & 61 & 5 & 0,08 & Hea \\
GPT-5.4 (Codex) & 47 & 1 & 0,02 & Väga hea \\
\end{tblr}
\end{table}

Tulemused näitavad nelja mudeli vahel märkimisväärset erinevust eesti keele ortograafilise täpsuse osas. GPT-5.4 Codex laienduse kaudu genereeritud töös tuvastati kogu 47-leheküljelise dokumendi peale vaid üks selge kirjaviga (\textit{vastusevormingun} õige \textit{vastusevormingu} asemel). Claude Opus 4.6 töös esinesid peamised vead pärisnimedes (\textit{Carnegei}, \textit{Californiya}, \textit{kanada} suurtähe puudumine) ning kahes tavasõnas (\textit{juhusiku}, \textit{stsenaariumm}), kokku viis juhtumit 61 lehekülje peale.

Claude Sonnet 4.6 töös leidus ortograafiliselt problemaatilisi sõnu rohkem: alapealkirjas kaks korda esinenud \textit{Masskiijutamine} (õige \textit{Masskirjutamine}), lühendite tabelis \textit{märkimisviiis}, ning veel mitu üksikut näpukat (\textit{kommertsiaalrsed}, \textit{rakenuste}, \textit{kasutusscenaariumide}, \textit{verstapostverstapost}). Osa vigadest asus eriti nähtaval kohal -- näiteks alapealkirjas või lühendite tabelis -- mis vähendab töö esmamuljet.

Kõige rohkem kirjavigu sisaldas Gemini 3.1 Pro genereeritud töö, kus 39 lehekülje kohta tuvastati 28 selget kirjaviga. Vead jagunesid nelja põhikategooriasse. Levinuima kategooria moodustasid puuduvate või vahetatud tähtedega tavasõnad (\textit{strukuure}, \textit{kordinaadid}, \textit{vektoorruumi}, \textit{konfiktide}, \textit{sepsifiiliste}, \textit{paering}, \textit{paringud}, \textit{uenit}, \textit{ekponentvõimendi}, \textit{hariduteadusele}). Teise rühma kuulusid valede tähtedega asendatud täpitähed (\textit{Öleksin} õige \textit{Oleksin} asemel, \textit{Üsaldan} õige \textit{Usaldan} asemel, \textit{Över-Reliance}, \textit{moita}). Kolmandana esinesid moodustatud mittesõnad (\textit{teoosiab}, \textit{esitlustuppu}, \textit{Tehisintsistendi}, \textit{algasestatud}). Neljanda kategooriana paistis silma ingliskeelse konventsiooniga termini \textit{biljonite} kasutamine eesti keeles sobiva \textit{miljardite} asemel. Lisaks esines üks pärisnime viga (\textit{Düning-Krugeri} õige \textit{Dunning-Krugeri} asemel).

Kirjavigade tihedus korreleerus hästi üldise kvaliteedihinnanguga: kõige madalamat kirjaviga lehekülje kohta näitas GPT-5.4 (agentse tööriista kaudu) ning Claude Opus 4.6. Gemini 3.1 Pro suurem vigade arv viitab, et Antigravity tööriist aitas küll mahtu ja struktuuri parandada, kuid ei suutnud täielikult kompenseerida mudeli nõrgemat eesti keele ortograafiat. See tulemus rõhutab, et automaatse korrektuuri tööriista (nt \texttt{aspell}, \texttt{hunspell}, JVKS) või täiendava keelelise toimetamise etapp on TI-ga genereeritud tekstide puhul vältimatu -- eriti mudelite korral, mille treeningandmetes on eesti keele osakaal väiksem.

\subsection{Mahu ja struktuuri analüüs} \label{subsec:mahu-analuus}

Lisaks kirjavigadele uuriti iga genereeritud töö mahtu ja sisemist struktuuri. Mahu hindamiseks loeti kokku kogu PDF-failist eraldatud sõnad (sealhulgas tekst, tabelite sisu, lisad ja viited) ning arvutati välja sõnade arv lehekülje kohta. Struktuuri hindamiseks loendati nummerdatud alapeatükkide, tabelite, jooniste ja koodinäidete arv kogu töös. Tulemused on esitatud tabelis~\ref{tab:mahuanaluus}.

\begin{table}[ht]
\centering
\caption{Genereeritud tööde maht ja struktuurielemendid}
\label{tab:mahuanaluus}
\begin{tblr}{
    colspec={l c c c c c c},
    row{1}={font=\bfseries},
    hlines,
    vlines,
    row{even}={bg=colorCellHighlight},
}
Mudel & Lk & Sõnu & Sõnu/lk & Alapt. & Tab. & Joon. \\
Gemini 3.1 Pro & 39 & 7\,230 & 185 & 17 & 2 & 0 \\
GPT-5.4 (Codex) & 47 & 9\,669 & 206 & 40 & 8 & 2 \\
Claude Sonnet 4.6 & 40 & 9\,739 & 243 & 24 & 12 & 2 \\
Claude Opus 4.6 & 61 & 10\,749 & 176 & 25 & 11 & 3 \\
\end{tblr}
\end{table}

Andmetest joonistuvad välja mitmed huvitavad mustrid. Kõige suurema mahu (10\,749 sõna) genereeris ootuspäraselt Claude Opus 4.6, mille väljund hõlmas ka kõige rohkem jooniseid (3) ning kompaktselt jaotatud sisu (176 sõna lehekülje kohta -- madalaim näitaja, mis tuleneb suuremast hulgast koodinäidetest, tabelitest ja diagrammidest). Claude Sonnet 4.6 paistis silma kõige tihedama tekstilise sisuga (243 sõna lehekülje kohta) ning sisaldas kõige rohkem tabeleid (12). GPT-5.4 (Codex) genereeris kõige granulaarsema struktuuriga töö -- 40 alapeatükki 47 lehekülje peale, mis viitab põhjalikule sisujaotusele. Gemini 3.1 Pro maht oli väikseim (7\,230 sõna), kõige vähemate alapeatükkidega (17), mis viitab pinnapealsemale käsitlusele.

Eraldi tasub märkida, et Gemini 3.1 Pro töös ei kasutatud ühtegi joonist (ainult tabeleid ja koodinäiteid), mis on bakalaureusetöö kontekstis ebatavaline -- enamik tehnilisi lõputöid sisaldab arhitektuuriskeeme või andmevoo diagramme. Sonnet 4.6 ja Gemini 3.1 Pro eelistasid suuremas mahus koodinäidete (Listings) lisamist (vastavalt 7 ja 5 koodinäidet), samas kui Opus 4.6 ja GPT-5.4 integreerisid koodinäited tihedamalt teksti sisse, ilma neid eraldi nummerdatud koodinäidete plokkidena vormistamata.

\subsection{Viidete kvantitatiivne analüüs} \label{subsec:viiteanaluus}

Akadeemilise töö üks olulisi kvaliteedinäitajaid on viidete kasutamise tihedus ja allikate hulk. Iga töö kohta loeti kokku kasutatud viidete koguarv (kirjanduse loendis esitatud unikaalsed allikad), tekstis tehtud tsitaatide arv (\texttt{[N]} tüüpi viidete esinemiskorrad enne kirjanduse loendit) ning arvutati tsitaatide tihedus lehekülje kohta. Tulemused on esitatud tabelis~\ref{tab:viited}.

\begin{table}[ht]
\centering
\caption{Genereeritud tööde viidete analüüs}
\label{tab:viited}
\begin{tblr}{
    width=\linewidth,
    colspec={X[1.4,l] X[1,c] X[1,c] X[1,c] X[1,c]},
    rows={font=\small},
    colsep=3pt,
    row{1}={font=\bfseries\small},
    hlines,
    vlines,
    row{even}={bg=colorCellHighlight},
}
Mudel & Unikaalseid viiteid & In-text tsitaate & Tsitaate lk kohta & Tsitaadi/viite suhe \\
Gemini 3.1 Pro & 9 & 13 & 0,33 & 1,44 \\
Claude Sonnet 4.6 & 17 & 23 & 0,58 & 1,35 \\
GPT-5.4 (Codex) & 22 & 28 & 0,60 & 1,27 \\
Claude Opus 4.6 & 35 & 45 & 0,74 & 1,29 \\
\end{tblr}
\end{table}

Tulemused näitavad, et Claude Opus 4.6 lisas kõige rohkem viiteid (35 unikaalset allikat ja 45 in-text tsitaati), saavutades ka kõrgeima tsitaaditiheduse 0,74 tsitaati lehekülje kohta. See vastab paroolihalduri valdkonna iseloomule, kus on vaja viidata krüptograafilistele standarditele (NIST, IETF RFC), turvauuringutele ja olemasolevatele lahendustele. GPT-5.4 (22 viidet) ja Claude Sonnet 4.6 (17 viidet) jäid keskmisesse vahemikku. Gemini 3.1 Pro paistis silma väikseima viidete arvuga (9 unikaalset allikat 39 lehekülje peale), mis on bakalaureusetöö kontekstis tagasihoidlik -- tüüpilistes arvutiteaduse bakalaureusetöödes on allikate arv sageli kõrgem kui kümme, eriti kirjanduse ülevaadet sisaldavate tööde puhul. See näitaja kinnitab üldist tähelepanekut, et Gemini 3.1 Pro genereeritud töö oli pigem mahukas-fabritseeritud kui allikatega põhjendatud.

Tsitaadi/viite suhe (in-text tsitaate jagatud unikaalsete allikate arvuga) näitab, mitu korda iga allikat keskmiselt korduvalt kasutati. Kõik mudelid jäid 1,27--1,44 vahemikku, mis viitab sellele, et enamikku allikaid kasutati tekstis vaid 1--2 korda. See on bakalaureusetööde puhul tüüpiline näitaja. Oluline on siiski rõhutada, et see analüüs ei hinda viidete \emph{sisulist} korrektsust ega kontrollitavust -- osa viidetest olid hallutsineeritud, eriti Sonnet 4.6 (ca 20\%) ja Opus 4.6 (ca 15\%) töödes.

\subsection{Anglitsismid ja eesti keele puhtus} \label{subsec:anglitsismid}

Eestikeelse akadeemilise teksti üks kvaliteedinäitajaid on tõlkimata jäetud või eesti keele konventsioonidest erinevate ingliskeelsete terminite vähene esinemine. Tartu Ülikooli ATI lõputööde juhend \cite{tuati2025juhend} rõhutab, et ``firmade, toodete ja programmide nimed ei ole terminid'' ning võõrkeelsed erialaterminid tuleks esimesel esinemisel tuua koos eestikeelse vaste ja kursiivis originaaliga. Selle põhjal on käesolevas analüüsis oluline eristada kolme erinevat tüüpi ingliskeelseid sõnu:
\begin{enumerate}
    \item \textbf{Vältimatud valdkonnaspetsiifilised tehnilised terminid} -- rahvusvaheliselt juurdunud standardi- või spetsifikatsiooniterminid, millel puudub ühtne ja üheselt arusaadav eestikeelne vaste (nt \textit{salt}, \textit{nonce}, \textit{hash}, \textit{token}, \textit{RAG}, \textit{transformer}). Need esinevad ka Eesti arvutiteaduse erialases kirjanduses ning nende säilitamine originaalkujul koos eestikeelse selgitusega on üldiselt aktsepteeritud.
    \item \textbf{Toote-, firma- ja tehnoloogianimed} -- nt \textit{Google}, \textit{Claude Code}, \textit{Cowork}, \textit{Antigravity}, \textit{OWASP}. Juhendi järgi ei ole need sisuliselt terminid ega anglitsismid, vaid pärisnimed ning ülakoma kasutamine käändelõppude eraldamiseks (\textit{Google'i}) on kooskõlas eesti keele konventsiooniga.
    \item \textbf{Välditavad inglismõjud} -- arvutiteaduse ingliskeelsed sõnad, millel on juurdunud eestikeelne vaste ja mis seetõttu on stiilivalikuna tõlgitavad (nt \textit{prompt}, \textit{pipeline}, \textit{benchmark}, \textit{over-reliance}, \textit{feature branch} ilma selgituseta). Need on kitsamas tähenduses ``anglitsismid'', mille hulk iseloomustab kõige otsesemalt eestikeelse stiili puhtust.
\end{enumerate}
Käesoleva analüüsi raames loendati esialgselt ülakomaga käändelõppudega vormistatud ingliskeelsed sõnad ja arvutiteaduse tehnilised terminid eristamata, ilma neid kolme kategooriasse süstemaatiliselt jaotamata; tulemused sellisel kujul on esitatud tabelis~\ref{tab:anglitsismid}. Analüüsist jäeti välja koodinäited ja kirjanduse loend.

\begin{table}[ht]
\centering
\caption{Genereeritud tööde anglitsismide analüüs (kategooriad eristamata)}
\label{tab:anglitsismid}
\begin{tblr}{
    width=\linewidth,
    colspec={X[1.5,l] X[1,c] X[1,c] X[0.9,c]},
    rows={font=\small},
    colsep=3pt,
    row{1}={font=\bfseries\small},
    hlines,
    vlines,
    row{even}={bg=colorCellHighlight},
}
Mudel+tööriist & Ülakomaga anglitsisme & CS-anglitsisme tekstis & Anglitsisme/lk \\
Claude Sonnet 4.6 + Cowork & 4 & 1 & 0,13 \\
GPT-5.4 + Codex & 9 & 11 & 0,43 \\
Claude Opus 4.6 + Cowork & 5 & 25 & 0,49 \\
Gemini 3.1 Pro + Antigravity & 7 & 22 & 0,74 \\
\end{tblr}
\end{table}

\textbf{Mõõdiku piirang ja tõlgendus.} Tabeli~\ref{tab:anglitsismid} ``Anglitsismide/lk'' näitaja koondab eeltoodud kolme kategooriat üheks arvuks ning on seetõttu tõlgenduslikult ebatäpne -- see karistab ühtmoodi nii välditavaid inglismõjusid kui ka vältimatuid standarditermineid ja pärisnimesid. Praktikas tähendab see, et valdkonnas, kus standardsed terminid on originaalkujul (krüptograafia, standardid, turvalisus), tõuseb mõõdik, kuigi tekst järgib aktsepteeritud akadeemilist konventsiooni. Claude Opus 4.6 + Cowork kombinatsiooni suhteliselt kõrge CS-anglitsismide arv (25) tuleneb suures osas just sellistest vältimatutest terminitest: \textit{salt}, \textit{nonce}, \textit{hash}, \textit{token}, \textit{AES-256-GCM}, \textit{Argon2id} on IETF/NIST standardites kinnistunud ega oma üheseid eestikeelseid vasteid, mida paroolihalduri kontekstis kasutatakse. Claude Sonnet 4.6 + Cowork TI-teema vältis seda probleemi osaliselt, kuna selle valdkonna terminoloogias on eesti keeles juba mitmeid kinnistunud vasteid. Gemini 3.1 Pro + Antigravity kasutab anglitsisme kõige rohkem (0,74/lk), sealhulgas paralleelselt nii eesti- kui ingliskeelseid vasteid (näiteks \textit{äärejuht (ingl. edge case)}). Täpsem kolmekategooriline loendus, mis eristaks välditavaid inglismõjusid vältimatutest terminitest ja pärisnimedest, annaks täpsema pildi ja on edasiste uurimuste oluline metoodiline täiendus; käesoleva töö tulemusi tuleks tõlgendada seda piirangut arvestades.

\subsection{Akadeemilise stiili indikaatorid} \label{subsec:akadeem-stiil}

Eestikeelses akadeemilises kirjutamises eelistatakse umbisikulist stiili (passiivi ja umbisikulist tegumoodi) ning välditakse esimese isiku ainsuse (\textit{ma, mina}) ja mitmuse (\textit{me, meie}) kasutamist. Iga genereeritud töö puhul loendati esimese isiku asesõnade (\textit{ma, mina, me, meie, meil, meid, minu}) esinemine peamises tekstis (välja arvatud litsents ja autorluse kinnitus). Tulemused on esitatud tabelis~\ref{tab:akadeem-stiil}.

\begin{table}[ht]
\centering
\caption{Esimese isiku asesõnade kasutus genereeritud töödes}
\label{tab:akadeem-stiil}
\begin{tblr}{
    colspec={l c c},
    row{1}={font=\bfseries},
    hlines,
    vlines,
    row{even}={bg=colorCellHighlight},
}
Mudel & Esimese isiku asesõnu & Hinnang \\
Claude Opus 4.6 & 0 & Suurepärane \\
GPT-5.4 (Codex) & 2 & Suurepärane \\
Claude Sonnet 4.6 & 2 & Suurepärane \\
Gemini 3.1 Pro & 4 & Hea \\
\end{tblr}
\end{table}

Kõik neli mudelit järgisid eestikeelset akadeemilist stiili väga hästi: esimese isiku asesõnu kasutati minimaalselt (0--4 esinemist kogu töö peale). Claude Opus 4.6 ei kasutanud ühtegi esimese isiku asesõna terve 61-leheküljelise töö peale, mis on ideaalne tulemus. Gemini 3.1 Pro neli esinemist on samuti vastuvõetav, kuid mõnevõrra kõrgem teistest. See näitab, et tänapäevased suured keelemudelid on eesti akadeemilise stiili konventsioonid hästi õppinud, mis on kooskõlas peatükis~\ref{subsec:tugevused} esitatud üldise hinnanguga TI mudelite akadeemilise stiili oskusele.

\subsection{Hinnangu stabiilsus ja subjektiivsus} \label{subsec:kooskola}

Kuna hindamise viis läbi üks hindaja (autor), ei olnud võimalik mõõta hindajatevahelist kooskõla. Subjektiivsuse vähendamiseks tugineti võimalikult palju kvantitatiivselt mõõdetavatele tunnustele (kirjavigade arv, lehekülgede ja sõnade arv, viidete arv, anglitsismide arv, esimese isiku asesõnade arv), mille korral ühe hindaja otsustusruum on piiratud ja mille tulemused on põhimõtteliselt reprodutseeritavad. Kvalitatiivsete hinnangute osas (näiteks teksti sisuline sügavus ja akadeemiline stiil) vaatas autor iga teksti läbi mitmel korral eri aegadel. Ühe hindajaga uurimisülesehituse piiranguid on käsitletud peatükis~\ref{arutelu}.
